\documentclass[11pt]{article}
\usepackage[letterpaper,margin=1in]{geometry}
\usepackage[T1]{fontenc}
\usepackage{lmodern,microtype,amsmath,amssymb,amsthm}
\usepackage[hidelinks]{hyperref}
\newtheorem{theorem}{Theorem}
\newtheorem{lemma}{Lemma}
\newtheorem{corollary}{Corollary}
\newcommand{\GL}{\operatorname{GL}}
\newcommand{\Sym}{\operatorname{Sym}}
\newcommand{\soc}{\operatorname{soc}}
\newcommand{\topmod}{\operatorname{top}}
\newcommand{\Span}{\operatorname{span}}
\title{Simple Young modules with nonprojective\\ polynomial injective hulls}
\author{Alec Kriebel\thanks{Independent researcher. ORCID:
\href{https://orcid.org/0009-0001-9320-500X}{0009-0001-9320-500X}.}}
\date{16 October 2026\quad (version 1.0)}
\begin{document}
\maketitle
\begin{abstract}
For every odd prime \(p\), the Young module labelled by
\(\lambda=(2p-2,1)\) in characteristic \(p\) is simple, whereas its corresponding
polynomial injective hull \(I_n(\lambda)\) is nonprojective for every \(n\ge2\).
At parameter \(q=1\), this gives a negative answer to the unrestricted
formulation of Miyachi's Conjecture~7 in the 2009 Oberwolfach report.
The proof combines classical representation theory with a labelled
rank-truncation argument. We include explicit proofs of the local calculations
and reproducible supporting checks. Historical priority remains unresolved;
no firstness claim is made.
\end{abstract}

\section{The printed question}
Let \(k\) be algebraically closed. In degree \(r\), write \(I_n(\nu)\) for the
injective hull of the simple polynomial \(\GL_n(k)\)-module \(L_n(\nu)\),
equivalently its injective hull for the Schur algebra \(S(n,r)\).
At stable rank \(N\ge r\), the Schur functor \(\mathsf F\) defines
\(Y^\nu=\mathsf F I_N(\nu)\), the Young module in the convention of
\cite[pp.~945--946]{M}. Here the parameter is \(q=1\), so the Hecke algebra is
\(k\mathfrak S_r\); in characteristic \(p\), its quantum characteristic is \(p\).

The first proposal in \cite[Conjecture~7, p.~946]{M} asks whether simplicity of
\(\topmod(Y^\nu)\) is equivalent to projectivity of \(I_n(\nu)\) for
\emph{some} \(n\ge\ell(\nu)\). The index-set criterion involving
\(\Lambda_n^i\) is presented there as a stronger proposal.
The characteristic-zero assumption in the subsequent subsection~0.5 concerns
a different method for simple Specht modules; it is not attached to the
displayed Conjecture~7.

\begin{theorem}\label{main}
Let \(p\) be any odd prime and let \(k\) be algebraically closed of
characteristic \(p\). Put \(r=2p-1\) and \(\lambda=(2p-2,1)\).
Then \(Y^\lambda\) is simple of dimension \(2p-2\), but \(I_n(\lambda)\)
is not projective for any integer \(n\ge2\).
\end{theorem}

Thus \(p=3,\lambda=(4,1)\) already refutes the broad printed proposal.
The theorem does not address a characteristic-zero variant with a nontrivial
root-of-unity parameter, nor the separate De Visscher--Donkin classification
of projective-injective polynomial modules. In particular, being row
\(p\)-regular is not being \(p\)-restricted (column \(p\)-regular), and no
projectivity inference from row regularity is used.

\section{Established structural inputs}
We use the following standard facts from \cite[Sections~2.2--2.4 and~3]{DD}.
Page locators here refer to the manuscript dated 12 January 2005
(printed pp.~7--12), whose repository cover is not counted.
Costandard reciprocity gives
\begin{equation}\label{reciprocity}
 (I_n(\nu):\nabla_n(\mu))=[\nabla_n(\mu):L_n(\nu)].
\end{equation}
A nonzero right side requires \(\mu\) to dominate \(\nu\).
An indecomposable polynomial injective is projective only if it is
contravariantly self-dual and an indecomposable tilting module.
Conversely, an injective tilting module is projective.
Contravariant duality fixes simple labels.

For \(N\ge n\), exact rank truncation takes \(I_N(\nu)\) to \(I_n(\nu)\)
and \(T_N(\mu)\) to \(T_n(\mu)\) when the respective labels have at most
\(n\) parts. We apply these labelled statements, not preservation of
projectivity for arbitrary modules under an idempotent functor.
Finally, \(\nabla_n(r)=\Sym^r(k^n)\) and
\(\nabla_2(a,b)=\det^b\otimes\Sym^{a-b}(k^2)\).
The tensor products of symmetric powers are injective, with injective
multiplicities given by the corresponding weight multiplicities of simples
\cite[Section~3, printed p.~12]{DD}.
None of these inputs assumes either conjecture. The stable Schur conventions
also appear in \cite[printed p.~3]{CP}.

\section{The Young module and its label}
Let \(V=k^r\) be the natural permutation module and put
\(A=\{(a_1,\ldots,a_r):\sum_i a_i=0\}\).
Since \(p\nmid r\), the constant line splits off and \(V=k\oplus A\).
For a nonzero submodule \(U\subseteq A\), choose \(0\ne v\in U\).
Its coordinates cannot all agree. If \(v_i\ne v_j\), then
\[
 v-(ij)v=(v_i-v_j)(e_i-e_j)\in U.
\]
Permuting the coordinates supplies all differences \(e_i-e_j\), which span
\(A\). Hence \(A\) is simple of dimension \(r-1\).
This is the classical natural permutation/augmentation calculation.

For clarity, we check the label in the source's \(\mathsf F I\) convention.
Take \(N\ge r\), \(E_N=k^N\) and
\(H=\Sym^{r-1}E_N\otimes E_N\).
Its injective multiplicities are the \((r-1,1,0,\ldots,0)\)-weight
multiplicities of simples. Dominance leaves only labels \((r)\) and
\(\lambda\), each with multiplicity one. For \(\lambda\) this is the highest
weight. For \(L_N(r)\), the weight occurs in the simple socle \(L_N(r)\) of
\(\nabla_N(r)=\Sym^r E_N\), which contains \(x_1^r\): the coefficient of \(t\) in
\((x_1+t x_2)^r\) is \(r x_1^{r-1}x_2\ne0\).
The ambient weight space is one-dimensional.
Thus
\[
 H\cong I_N(r)\oplus I_N(\lambda).
\]
Reciprocity at the maximal label gives \(I_N(r)=\Sym^r E_N\).
Multiplication \(m:H\to\Sym^r E_N\) splits: the polarization map
\[
 \delta(f)=\sum_i \frac{\partial f}{\partial x_i}\otimes x_i
 \quad\text{satisfies}\quad m\delta=r\,\mathrm{id}.
\]
Cancellation of indecomposable summands yields \(\ker m\cong I_N(\lambda)\).
On multilinear weight spaces, \(H\) has basis
\((\prod_{j\ne i}x_j)\otimes x_i\) for \(1\le i\le r\), and \(m\) is
coordinate summation onto \(k\prod_jx_j\). Exactness of \(\mathsf F\)
therefore gives \(\mathsf F I_N(\lambda)\cong A\), as required.
The stable Schur functor is never applied at rank \(2<r\).

\section{The rank-two obstruction}
Write \(E=k^2\), with basis \(x,y\). In a fixed symmetric power,
the monomials \(v_i=x^{m-i}y^i\) have distinct algebraic-torus weights.
An invariant subspace contains individual weight components.
It also contains every coefficient in \(t\) of a unipotent translate:
Vandermonde inversion over the infinite field \(k\) extracts those
coefficients from evaluations. This is a statement about rational
\(\GL_2(k)\)-modules, not just the finite group \(\GL_2(\mathbb F_p)\).

\begin{lemma}\label{local}
The module \(\Sym^{2p-1}E\) is simple. The module
\(B=\Sym^{2p-3}E\) is a nonsplit extension of two distinct simple modules:
its socle has highest weight \((2p-3,0)\) and dimension \(2p-4\);
its top has highest weight \((p-1,p-2)\) and dimension \(2\).
\end{lemma}
\begin{proof}
From any monomial, the substitution \(y\mapsto y+t x\) extracts the
highest monomial. Translating that highest monomial in the opposite
direction extracts the binomial-coefficient multiples of all monomials.
For degree \(2p-1\), every coefficient is nonzero, since in characteristic \(p\)
\[
 (1+t)^{2p-1}=(1+t^p)(1+t)^{p-1},
\]
and the coefficients of \((1+t)^{p-1}\) are alternating units.
This proves the first assertion.

For \(m=2p-3\), let
\[
 W=\Span\{v_i:0\le i\le p-3\text{ or }p\le i\le2p-3\}\subset B.
\]
It is the image of the injective equivariant multiplication
\(E^{(1)}\otimes\Sym^{p-3}E\to B\), where \(E^{(1)}\) denotes Frobenius
twist, via \(x^{(1)}\otimes f\mapsto x^p f\) and
\(y^{(1)}\otimes f\mapsto y^p f\). The two monomial ranges are disjoint. The same highest-monomial
argument, now using
\[
 (1+t)^{2p-3}=(1+t^p)(1+t)^{p-3},
\]
shows that every nonzero submodule of \(W\) contains all its displayed
monomials. Hence \(W\) is simple with the stated highest weight and dimension.

The quotient has the two weight lines represented by \(v_{p-2},v_{p-1}\).
The opposite unipotents connect them with coefficient \(p-1\ne0\), so the
quotient is simple of highest weight \((p-1,p-2)\).
Any invariant complement to \(W\) must be a sum of torus weight spaces,
hence must equal \(\Span(v_{p-2},v_{p-1})\).
But the translate of \(v_{p-2}\) under \(y\mapsto y+t x\) has the
nonzero coefficient \(x^{2p-3}\in W\) at \(t^{p-2}\).
There is no complement. A nonsplit extension of two distinct simples
has precisely the asserted socle and top.
\end{proof}

These rank-two structures are established representation theory;
for published formulas yielding them by duality and determinant twist, see
\cite[Proposition~3.1, classical analogue, and Theorem~3.3]{EL}.
The explicit proof makes their use here checkable.

Only \((r)\) strictly dominates \(\lambda=(r-1,1)\).
By Lemma~\ref{local}, \(\nabla_2(r)=\Sym^r E\) is simple and does not
contain \(L_2(\lambda)\). Equation~\eqref{reciprocity} therefore gives
\[
 I_2(\lambda)=\nabla_2(\lambda)=\det\otimes B.
\]
Its socle and top are \(L_2(2p-2,1)\) and \(L_2(p,p-1)\), respectively.
Their highest weights differ, including when \(p=3\) and both dimensions
are \(2\). Thus \(I_2(\lambda)\) is not contravariantly self-dual and
is not projective.

\section{Excluding every permitted rank}
Suppose \(I_N(\lambda)\) were projective for some \(N\ge2\).
It would be \(T_N(\mu)\) for a partition \(\mu\) of \(r\).
Since \(L_N(\lambda)\) is a composition factor, \(\mu\) dominates \(\lambda\).
Consequently \(\mu=(r)\) or \(\lambda\); both labels have at most two parts.
Labelled rank truncation gives
\[
 I_2(\lambda)\cong T_2(\mu).
\]
This would be an injective tilting module, hence projective, contradicting
the rank-two calculation. This proves Theorem~\ref{main} for every \(N\ge2\).

\begin{corollary}
The stronger simple-top/\(\Lambda_n^i\) criterion in the same printed
Conjecture~7 is false in this setting.
\end{corollary}
\begin{proof}
Membership in the constructed \(\Lambda_n^i\) family is already a sufficient
condition for \(I_n(\lambda)\) to be projective \cite{DD,M}.
Theorem~\ref{main} excludes membership at every allowed rank.
No completeness assertion about that family is needed.
\end{proof}

\section{Verification, attribution and status}
The companion archive preserves the frozen candidate, independent
mathematical and priority reviews, source identities and reproducible exact
checks. Formal polynomial-coefficient calculations in characteristics
\(3\) and \(5\) verify the local constituents and nonsplitting.
A separate tensor-orbit computation constructs all \(56\) basis elements
of \(S(2,5)\), checks its \(3136\) products and independently finds
\(\dim P_2(3,2)=6\), whereas \(\dim I_2(4,1)=4\).
That optional computation uses NumPy; the core checks use Python's standard
library. Finite computations support the argument, and do not prove its
all-prime or all-rank assertions.

The Young/augmentation calculation, rank-two structure and Schur-algebra
inputs are classical. The note's purpose is their explicit assembly into
a counterexample to the broad printed 2009 proposal.
An extensive but bounded literature audit on 16 October 2026 did not
verify an earlier complete bridge refutation. Priority remains unresolved:
some older two-row papers, tables and later endomorphism-algebra papers
were not available in full, and citation coverage is incomplete.
No claim of firstness, certified novelty, or current worldwide open status
is made. The source archive records exact read scopes and gaps.

AI tools were used extensively in developing and checking the argument,
literature investigation, computations, drafting, reproduction and
adversarial review. This is an unrefereed preprint without external human
peer review or proof-assistant certification. The note and authored
companion material are licensed under CC BY 4.0.

\begin{thebibliography}{9}
\raggedright
\bibitem{M}
H.~Miyachi,
Dipper's hypothesis and self-injective endomorphism rings,
in \emph{Representations of Finite Groups},
Oberwolfach Reports \textbf{6} (2009), report~17, pp.~944--947.
\href{https://doi.org/10.4171/OWR/2009/17}{doi:10.4171/OWR/2009/17}.

\bibitem{DD}
M.~De Visscher and S.~Donkin,
On projective and injective polynomial modules,
\emph{Math. Z.} \textbf{251} (2005), 333--358.
\href{https://doi.org/10.1007/s00209-005-0805-x}{doi:10.1007/s00209-005-0805-x}.
Specific manuscript locators refer to the version dated 12 January 2005:
\url{https://openaccess.city.ac.uk/id/eprint/969/1/prinjjan12.pdf}.

\bibitem{CP}
A.~Cox and A.~Parker,
Homomorphisms and higher extensions for Schur algebras and symmetric groups,
\emph{J. Algebra Appl.} \textbf{4} (2005), 645--670.
\href{https://doi.org/10.1142/S0219498805001460}{doi:10.1142/S0219498805001460}.
The cited printed p.~3 is in the institutional accepted version:
\url{https://openaccess.city.ac.uk/id/eprint/373/2/coxparker2final.pdf}.

\bibitem{EL}
K.~Erdmann and S.~Law,
Torsion pairs and Ringel duality for Schur algebras,
\href{https://arxiv.org/abs/2104.05317v1}{arXiv:2104.05317v1} (12 April 2021),
pp.~5--6.
\end{thebibliography}
\end{document}
